% Layout extracted from Kupzyk report at ca1d062; retain a descriptive title.
\documentclass[11pt]{article}
\usepackage[letterpaper,margin=0.82in]{geometry}
\usepackage{booktabs,longtable,array,graphicx,amsmath,xcolor,hyperref,xurl}
\usepackage[authoryear,round]{natbib}
\hypersetup{colorlinks=true,urlcolor=blue,linkcolor=blue,citecolor=blue}
\renewcommand{\arraystretch}{1.25}
\setlength{\emergencystretch}{1em}
\setlength{\parskip}{5pt}\setlength{\parindent}{0pt}
% Use scalable fonts; inspect pdffonts output after building. Avoid adding T1
% fontenc without an installed scalable font family. Add title/date in manuscript.

\usepackage{float}
\ExplSyntaxOn
\cs_new_protected:Npn \koziol_showpath:n #1
  {
    \tl_set:Nn \l_tmpa_tl {#1}
    \tl_replace_all:Nnn \l_tmpa_tl { / } { /\allowbreak }
    \group_begin:
      \ttfamily
      \tl_use:N \l_tmpa_tl
    \group_end:
  }
\cs_new_eq:NN \showsrcpath \koziol_showpath:n
\ExplSyntaxOff
\NewDocumentCommand{\src}{v}{%
  \href{https://github.com/alectionary/vibe-math/blob/master/formalization/koziol-weighted-multilevel/#1}{\showsrcpath{#1}}%
}

\title{Weighted single-level versus multilevel:\\
a replication of Koziol (2015)}
\date{3 November 2026}
\begin{document}
\maketitle

\begin{abstract}
We replicated the headline result of Koziol (2015) on the cells that carry it:
weighted single-level estimators tended to outperform weighted multilevel
estimators, and our contrasts are 0.152 against the printed 0.151 for the
normal intercept (998 of 1,000 replications), $-0.173$ against $-0.170$ for
the Bernoulli threshold (100 replications; $\tau=-\gamma$ in
$P(Y=1)=\Phi(\gamma+\cdots)$), and 0.033 against 0.033 for the
random-coefficient approximate-weight minus true-weight slope (40
replications) \citep{koziol2015}.
Not all 208 settings were rerun: coverage of the weighted multilevel intercept
at average cluster size 20 and intraclass correlation .05 is 0.252 against the
printed 0.282, the factors $G/(G-1)$ and $N/(N-1)$ move our coverage by about
0.001, and the Monte Carlo standard error of that coverage is 0.014.
The source software is Mplus~7.
The tolerance on a difference of two three-decimal biases is 0.001, and a
positive difference means our Monte Carlo mean is larger than the printed
value.
\end{abstract}

\begin{table}[ht!]
\centering
\footnotesize
\setlength{\tabcolsep}{3.5pt}
\renewcommand{\arraystretch}{1.15}
\caption{Signed contrasts at the cells that carry the headline.
Published values are the printed contrasts in Tables C.2, C.4, and E.2 of
Koziol (2015).
Our values are Monte Carlo means.
Rows 1, 2, and 4 are in outcome units.
Row 3 is in probit threshold units.
Row 5 is a proportion.
The tolerance on a difference of two three-decimal biases is 0.001.
Magnitude uses the unrounded mean: a gap of at most 0.001 is same.}
\label{tab:headline}
\begin{tabular}{@{}p{2.05in}ccccc@{}}
\toprule
Estimand & Published & Ours & Effect magnitude & Direction & Verdict \\
\midrule
Normal wtm minus ws intercept & 0.151 & 0.152 & larger & same direction & mismatch \\
Approx.\ minus true intercept & 0.157 & 0.156 & same & same direction & close \\
Bernoulli wtm minus ws threshold & $-0.170$ & $-0.173$ & larger & same direction & mismatch \\
RC slope, approx.\ minus true & 0.033 & 0.033 & same & same direction & match \\
Normal ws minus wtm coverage & 0.945 & 0.945 & same & same direction & match \\
\bottomrule
\end{tabular}

\vspace{0.4em}
\begin{minipage}{0.98\textwidth}
\footnotesize
The source software is Mplus~7.
The tolerance on a difference of two three-decimal biases is 0.001.
A positive difference means our Monte Carlo mean is larger than the printed value.
A positive estimand means the first estimator named in that row exceeds the second.
Descriptive point estimates; not equivalence.
\end{minipage}
\end{table}

\section{Introduction}

We replicated the headline result of Koziol (2015) on the cells that carry it.
Under informative sampling, weighted single-level estimators tended to
outperform weighted multilevel estimators.
Our normal intercept contrast is 0.152 against the printed 0.151, on 998 of
1,000 replications.
Our Bernoulli threshold contrast is $-0.173$ against $-0.170$, on 100
replications, with $\tau=-\gamma$ in $P(Y=1)=\Phi(\gamma+\cdots)$.
Our random-coefficient approximate-weight minus true-weight slope contrast is
0.033 against 0.033, on 40 replications.
Not all 208 settings were rerun.
At average cluster size 20 and intraclass correlation .05, weighted multilevel
coverage is 0.252 against the printed 0.282.
The factors $G/(G-1)$ and $N/(N-1)$ move that coverage by about 0.001, and
the Monte Carlo standard error of the coverage is 0.014.
The source software is Mplus~7 \citep{muthen2012}.
The tolerance on a difference of two three-decimal biases is 0.001.
A positive difference means our Monte Carlo mean is larger than the printed
value.
These comparisons are descriptive point estimates, not equivalence tests.

Natalie A.\ Koziol concluded that, when selection probabilities are unequal
and informative, a weighted single-level estimator with sandwich standard
errors recovers the intercept more accurately than a weighted multilevel
pseudo-maximum likelihood estimator, and that approximate cluster weights
leave more bias than the true multilevel weights
\citep{koziol2015}.
The dissertation is \emph{A Comparison of Population-Averaged and
Cluster-Specific Approaches in the Context of Unequal Probabilities of
Selection} (University of Nebraska--Lincoln, 2015), chaired by James A.\
Bovaird, and archived at
\url{https://digitalcommons.unl.edu/cehsdiss/229/}.
The three research questions ask how population-averaged and cluster-specific
estimators behave under informative sampling, how their interval estimates
cover, and whether an approximate cluster weight can stand in for the true
multilevel weight.
Table~1 answers those questions on the cells we reran.
A population-averaged fit is a single-level model with a sandwich variance.
A cluster-specific fit is a multilevel model with a random intercept, or with
a random intercept and a random slope.
A sample is informative when the selection probability depends on the cluster
effect that also enters the outcome, so an unweighted fit estimates a
different intercept from the census model.

Five estimators appear in the tables.
The labels below are hers.
Unweighted multilevel (\texttt{um.ccs}) and unweighted single-level
(\texttt{us.ics}) ignore the weights.
Weighted single-level (\texttt{ws.ics}) uses the selection weights and a
cluster sandwich \citep{binder1983,lumley2004}.
Weighted true multilevel (\texttt{wtm.ccs}) uses the multilevel weights inside
a pseudo-likelihood \citep{pfeffermann1998,asparouhov2006,rabehesketh2006}.
Weighted approximate multilevel (\texttt{wam.ccs}) replaces the cluster weight
by $M$ divided by the realized number of clusters, and divides the conditional
weight by that constant, the approximation of \citet{kovacevic2003}.
$M$ is the number of clusters in the population.
The population-averaged estimating equation is the same class of model that
geepack implements for generalized estimating equations
\citep{hojsgaard2006}.
The single-level fits in this note use the survey package, because the
design is a probability sample with a cluster sandwich
\citep{lumley2024}.

\section{What do the three designs show?}

The question is what the normal random-intercept grid, the Bernoulli
random-intercept cells, and the random-coefficient cell show.
The answer is that the estimator order in her tables reappears in each design
we reran, and that one coverage rate at the larger cluster size does not.

\subsection{Normal random intercept}

The question is whether weighted multilevel intercept bias exceeds weighted
single-level intercept bias when the outcome is normal.
The answer is yes, on the primary cell and on the rest of the eight-cell grid.

The primary cell is informative sampling, intraclass correlation .05, and
average cluster size 5.
Koziol reported the contrast of \texttt{wtm.ccs} minus \texttt{ws.ics}
intercept bias as 0.151 in Table C.2.
Our contrast on 998 retained replications is 0.152.
The paired Monte Carlo standard error is 0.000416.
The unrounded gap is 0.00103, which is a mismatch under the 0.001 rule.
If 0.151 is treated as exact, the gap is slightly larger than two of our
standard errors.
The order is unchanged: weighted multilevel bias is the larger positive number.
In the same cell the approximate-weight contrast, \texttt{wam.ccs} minus
\texttt{wtm.ccs}, is 0.156 against the printed 0.157, with paired standard
error 0.000515.
Both contrasts are in outcome units.
The cell means are stored in \src{simulation/results/summary.csv}, and the
replicate file is \src{simulation/results/cell_01.csv}.

Our intercept biases on that cell are 0.682 unweighted, 0.309 with
approximate weights, 0.153 with true multilevel weights, and 0.001 with the
weighted single-level estimator.
The printed true-multilevel bias is 0.151, so the raw multilevel mean, 0.153,
misses the printed thousandth by 0.002, and the contrast with the
single-level estimator is a mismatch at the same rule.
Weighted single-level coverage of the intercept is 0.945, and weighted
multilevel coverage is 0, so the coverage contrast is 0.945 against the
printed 0.945 in Table E.2.
A coverage near .95 on 998 draws has a binomial standard error near 0.007.
She reported that the weighted single-level intervals sit near the nominal
rate and the weighted multilevel intervals do not.
Our run shows the same split.
Two of the 1,000 requested replications were dropped for nonconvergence.

The other informative cells keep the same order.
At intraclass correlation .25 and average cluster size 5, our true-multilevel
intercept bias is 0.163 and the weighted single-level bias is 0.003.
At average cluster size 20 the true-multilevel bias falls to 0.045
(intraclass correlation .05) and 0.047 (intraclass correlation .25), while the
weighted single-level bias stays near 0.002.
Approximate weights remain worse than true weights in every informative cell.
Under non-informative sampling the intercept biases are small.
Across the four non-informative cells and five estimators, her printed
intercept biases lie between $-0.002$ and 0.001, and ours lie between
$-0.002$ and 0.002.
The largest absolute bias is 0.002 in both sources (Table C.1).
Relative bias is undefined at a zero intercept, so the comparison is the
absolute bias.

The realized primary sample has mean 967.7 clusters and 4,884 students, and
the mean cluster weight is 49.5.
Table 4.1 reports a mean of 964 clusters and a mean cluster weight of about
50.
A difference of about four clusters on a mean near 964 is under one percent.
Cluster counts are means of Poisson samples drawn with separate seeds, so
exact equality is not expected.
The design check supports the bias comparison.
It is not itself a performance result.

\subsection{Bernoulli random intercept}

The question is whether the weighted multilevel threshold is more negative
than the weighted single-level threshold when the outcome is Bernoulli.
The answer is yes, at intraclass correlation .05 and average cluster size 5,
on 100 replications.

The outcome model used here is $P(Y=1)=\Phi(\gamma+x\beta+u)$.
Table C.2's intercept column matches the threshold $\tau=-\gamma$, not the
latent intercept $\gamma$.
At average cluster size 5 our threshold means are $-0.715$ unweighted against
the printed $-0.712$, $-0.002$ weighted single-level against 0.000, $-0.174$
true-weight multilevel against $-0.170$, and $-0.364$ approximate-weight
against $-0.361$.
The contrast of true-weight multilevel minus weighted single-level is
$-0.173$ against the printed $-0.170$.
The paired Monte Carlo standard error on 100 replications is 0.00160.
The gap of 0.003 is larger than the 0.001 tolerance and inside two of our
standard errors (0.0032).
The sign is the same, and the four estimator means sit in the same order as
Table C.2.
These means are in \src{simulation/results/bern_inf_05_5.csv}.
Her Bernoulli cells used 1,000 replications.
Ours use 100.

At average cluster size 20 the same pattern is smaller, as she reported.
Our true-weight threshold mean is $-0.042$ against the printed $-0.043$, and
the approximate-weight mean is $-0.207$ against $-0.208$
(\src{simulation/results/bern_inf_05_20.csv}).
The unweighted threshold mean is $-0.713$.
The weighted single-level mean is 0.001.
Only these two Bernoulli cells, both at intraclass correlation .05, were
rerun.

\subsection{Normal random coefficient}

The question is whether approximate cluster weights bias the level-1 slope
when the true multilevel weights do not.
The answer is yes, on the one informative cell we reran.

The cell is a normal outcome, informative sampling, intraclass correlation
.05, and average cluster size 5.
Table C.4 prints the contrast of approximate-weight slope bias minus
true-weight slope bias as 0.033.
Our contrast on 40 replications is 0.033.
The paired Monte Carlo standard error is 0.00190.
The replicate file is \src{simulation/results/rc_inf_05_5.csv}.
All 40 replications converged.
The slope bias subtracts the common true slope stored on each row.

The component biases, ours against the printed value, are these.
The unweighted intercept is 0.687 against 0.684, and the unweighted slope is
0.035 against 0.034.
The weighted single-level intercept is 0.012 against 0.002, with Monte Carlo
standard error 0.004, and the slope is 0.002 against 0.002.
The true-weight multilevel intercept is 0.180 against 0.174, and the slope is
0.004 against 0.003.
The approximate-weight intercept is 0.337 against 0.333, and the slope is
0.037 against 0.036.
The slope contrast is the result that carries this design.
The weighted single-level intercept differs from the printed 0.002 by 0.010,
about two and a half of our standard errors, on only 40 replications.
That intercept comparison is too noisy to settle, and it is not the contrast
in Table~1.

\section{What did we do, and what did we take from Koziol?}

The question is which parts of the numerical work are ours and which parts are
taken from the dissertation.
The answer is that the design, the target parameters, and the printed tables
are hers, and the Monte Carlo means in this note are ours.

Koziol specified a finite population, Poisson sampling at the cluster and
student levels, normal and Bernoulli outcomes, random-intercept and
random-coefficient structures, intraclass correlations .05 and .25, and
average cluster sizes 5 and 20.
Chapter III and Appendix A give the operational weight formulas and a fragment
of the R generator.
Page 76 names Mplus version 7 for estimation, with \texttt{TYPE=COMPLEX} or
\texttt{TWOLEVEL}, \texttt{ESTIMATOR=MLR}, \texttt{LINK=PROBIT},
\texttt{WTSCALE=CLUSTER}, and \texttt{BWTSCALE=UNSCALED}
\citep{muthen2012}.
Tables C.2, C.4, and E.2 supply the published biases and coverages used above.
Table 4.1 supplies the published sample-size moments.
The 2017 article in \emph{Multivariate Behavioral Research} repeats the
comparison \citep{koziol2017}.
Its Figshare supplement repeats tables and does not contain Mplus or R
programs.
The article PDF was not re-read for this note.
Every published number below is marked as taken from those tables.
Every Monte Carlo mean is marked as our simulation.
Choices that her deposited programs would have fixed, and that we set in order
to run, are marked as assumptions in the next section.

Our reconstruction is in R 4.3.3 \citep{r2024},
\url{https://www.R-project.org/}.
The dissertation generates the populations in R with MASS.
Our generator is \src{simulation/generate.R}.
Weighted single-level fits and their cluster-robust standard errors use survey
4.2.1 \citep{lumley2004,lumley2024},
\url{https://CRAN.R-project.org/package=survey}.
Unweighted Gaussian multilevel fits, and the unweighted probit start, use
lme4 1.1.35.1 \citep{bates2015},
\url{https://CRAN.R-project.org/package=lme4}.
The random-slope pseudo-likelihood, and a one-replicate check of the normal
multilevel fit, use WeMix 4.0.3 \citep{bailey2023},
\url{https://CRAN.R-project.org/package=WeMix}.
On that production replicate, WeMix and the analytic pseudo-likelihood agree.
Numerical derivatives for the pseudo-likelihood sandwich use numDeriv
2016.8-1.1 \citep{gilbert2019},
\url{https://CRAN.R-project.org/package=numDeriv}.
The Bernoulli multilevel estimator is a Gauss--Hermite pseudo-likelihood with
15 nodes, started at the unweighted \texttt{glmer} fit and optimized with
\texttt{nlminb}.
The run record is \src{simulation/results/RUN.txt}.

The normal random-intercept factorial was run at 1,000 requested replications
per cell.
The two Bernoulli cells were run at 100 replications.
The random-coefficient cell was run at 40 replications.
Mplus 7 was not available for a refit of the same samples, and the original
input files were not found.

\section{What assumptions did we supply?}

The question is which choices we supplied because the estimation programs were
not deposited.
The answer is the seven choices in this section.
A setting she states in the dissertation stays hers.

\subsection{Threshold sign}

The question is which Bernoulli sign we compare with Table C.2.
The answer is the threshold $\tau=-\gamma$, where $\gamma$ is the intercept
in $P(Y=1)=\Phi(\gamma+x\beta+u)$.

The latent intercept has the opposite sign from the column she printed.
The threshold matches that column, including the unweighted mean near $-0.71$
and the approximate-weight mean near $-0.36$.
This is a parameterization that matches the table.
It is not a finding that the table is wrong.

\subsection{Method 2 scaling}

The question is how the multilevel weights are scaled before the
pseudo-likelihood.
The answer is the cluster scaling she names as \texttt{WTSCALE=CLUSTER},
which this reconstruction treats as Method 2, with between-cluster weights
left unscaled (\texttt{BWTSCALE=UNSCALED}).

Asparouhov describes scaling choices for multilevel sampling weights, and the
Mplus 7 user's guide names the option on the page she cites
\citep{asparouhov2006,muthen2012,rabehesketh2006}.
We did not refit the grid under the other scale options.
A different scale would change the weighted multilevel estimator, so this
choice is part of the comparison rather than a harmless default.

\subsection{Normal intervals at 1.96}

The question is which critical value defines coverage.
The answer is 1.96.

Coverage is the proportion of replications in which the Wald interval
$\hat\theta\pm 1.96\,\mathrm{SE}$ contains the target.
We did not use a $t$ critical value and we did not apply a further
degrees-of-freedom correction on top of the sandwich.
Table E.2 is compared with that definition.
The single-level standard errors are the survey cluster sandwich.
The weighted multilevel standard errors are the observed pseudo-likelihood
sandwich.

\subsection{Separate seeds}

The question is whether our samples are her samples.
The answer is no.

Her seed stream was not deposited.
Each cell in this reconstruction uses its own seed sequence.
Monte Carlo means are the objects compared with the printed tables.
A single replicate is not expected to match a replicate she printed, and she
did not print replicate-level estimates.
The cluster-count gap of about four clusters is the sort of difference
separate seeds produce.

\subsection{Model-based standard errors in lme4}

The question is where the unweighted Gaussian multilevel standard errors come
from.
The answer is the model-based variance from lme4, not an Mplus MLR sandwich.

The weighted fits use a sandwich, as her MLR columns do.
The unweighted multilevel coverage in our normal grid is therefore not a
reconstruction of \texttt{ESTIMATOR=MLR}.
The headline coverage contrast uses the weighted single-level sandwich and the
weighted multilevel sandwich, so this assumption does not enter Table~1.
It does limit any side comparison of unweighted multilevel coverage with
Table E.2.

\subsection{Appendix A variances}

The question is which variance components generate the stated intraclass
correlation.
The answer is the pair printed in the Appendix A generator fragment.

The fragment fixes the scale that turns intraclass correlations .05 and .25
into a level-2 variance and a residual variance.
We treated that fragment as the complete normal generator and did not estimate
the variances from her tables.
If the fragment we transcribed omits a scale factor, every normal bias would
move by that factor.
The Bernoulli threshold is on the probit scale and does not use those
Gaussian variances.

\subsection{Cluster-aligned random slope}

The question is how the random slope enters the linear predictor.
The answer is $u_{1j} x_{ij}$, with the slope draw aligned to cluster $j$.

WeMix then fits the weighted random-slope pseudo-likelihood under Method 2
scaling \citep{bailey2023}.
A slope vector that is not aligned to the cluster estimates a different
covariance and a different level-1 coefficient.
The match of the slope contrast to Table C.4 depends on this alignment.

\section{Which ambiguities were resolved?}

The question is which implementation choices were checked against an
alternative and then fixed.
The answer is the four checks below.
Each one changes which number is allowed into the tables.

\subsection{The two probit optimizers agree}

The ambiguity was whether the Gauss--Hermite pseudo-likelihood had settled on
a local mode.
The check was a second optimization with BFGS, from the same start, on every
10th replication of the weighted multilevel probit.
The finding is that the largest absolute gap on the latent intercept was
under 0.0003 in both Bernoulli cells
(\src{simulation/results/bern_inf_05_5.csv},
\src{simulation/results/bern_inf_05_20.csv}).
All 100 replications in each cell converged.
The implication is that the threshold contrast is not an optimizer artifact.
The reported means use \texttt{nlminb}.

\subsection{Observation weights in glmmTMB do not match Table C.2}

The ambiguity was whether a weighted probit mixed model with observation
weights is the same estimator as the multilevel pseudo-likelihood.
The check was glmmTMB with observation weights on the same Bernoulli samples,
\url{https://CRAN.R-project.org/package=glmmTMB}.
The finding is that the threshold mean at average cluster size 5 is about
$-0.974$, against the printed true-weight mean of $-0.170$.
At average cluster size 20 the glmmTMB threshold mean is about $-0.247$,
against the printed $-0.043$.
The implication is that this likelihood is a different estimator.
It is not the contrast in Table~1, and it is not evidence against Table C.2.

\subsection{GLMMadaptive weights are per cluster}

The ambiguity was whether a mixed-model package that accepts a weight argument
applies that weight once per observation.
The check was the weight behavior of GLMMadaptive,
\url{https://CRAN.R-project.org/package=GLMMadaptive}.
The finding is that the package applies the weight once per cluster, not once
per observation.
The implication is that those fits are not the observation-level
pseudo-likelihood behind Table C.2.
They were not used as the reported Bernoulli means.

\subsection{A recycled slope vector was discarded}

The ambiguity was whether the random slope was indexed by cluster.
The check was a first implementation that recycled a slope vector down the
rows, followed by a refit in which the slope draw is constant inside the
cluster, $u_{1j}x_{ij}$.
The finding is that the recycled vector does not estimate the level-1 slope
she printed.
It was discarded.
After the index was corrected, the approximate-weight minus true-weight slope
contrast on 40 replications is 0.033 against the printed 0.033
(\src{simulation/results/rc_inf_05_5.csv}).
A separate 20-replication check compared a raw covariate with a centered
covariate in the unweighted \texttt{lmer} fit.
The reported random-coefficient results use the cluster-aligned slope.

\section{How do the numbers compare?}

The question is which printed numbers match, which are close, and which miss,
once scope is separated from the third decimal.
The answer is that the slope contrast and the primary coverage contrast match,
the approximate-weight intercept contrast is close, and the normal intercept
contrast, the Bernoulli threshold contrast, and the cluster-size-20 coverage
are mismatches at the thousandth.
Scope is limited on every row: the samples are a reconstruction in R, not
Mplus~7, and not her seed stream.

Numerical agreement uses three labels at the printed precision.
Match means our value rounds half up to the printed value.
Close means it does not, and the absolute gap is at most one unit in the last
printed place.
Mismatch means the gap is larger than that.
Scope is limited when the estimator and the scientific target are the ones she
defined but the software, the seeds, or the replication count differ.
Limited scope does not change a numerical match.
The page-1 tolerance of 0.001 is the sum of two half-thousandth rounding
bounds, one for each bias that enters a contrast.
That tolerance is a descriptive rule for Table~1.
It is not an equivalence margin, and it is not the Monte Carlo standard error.

\begin{table}[H]
\centering
\footnotesize
\setlength{\tabcolsep}{4pt}
\renewcommand{\arraystretch}{1.3}
\caption{Numerical agreement and reconstruction scope.
Published values are taken from the dissertation tables named in the text.
Our values are from this simulation.
Precision is the printed unit used for the numerical label.}
\label{tab:compare}
\begin{tabular}{@{}p{2.15in}rrccp{0.72in}@{}}
\toprule
Comparison & Published & Ours & Precision & Numerical & Scope \\
\midrule
Normal wtm minus ws intercept & 0.151 & 0.152 & 0.001 & mismatch & limited \\
Normal approx.\ minus true intercept & 0.157 & 0.156 & 0.001 & close & limited \\
Normal ws minus wtm coverage & 0.945 & 0.945 & 0.001 & match & limited \\
Bernoulli wtm minus ws threshold & $-0.170$ & $-0.173$ & 0.001 & mismatch & limited \\
Random-coefficient slope contrast & 0.033 & 0.033 & 0.001 & match & limited \\
Coverage at cluster size 20 & 0.282 & 0.252 & 0.001 & mismatch & limited \\
Raw normal wtm intercept bias & 0.151 & 0.153 & 0.001 & mismatch & limited \\
Largest non-informative absolute bias & 0.002 & 0.002 & 0.001 & match & limited \\
Mean cluster count, primary cell & 964 & 967.7 & 1 & mismatch & limited \\
\bottomrule
\end{tabular}
\end{table}

The normal intercept contrast is a mismatch because the unrounded gap is
0.00103, which is larger than 0.001, and 0.152 does not round to 0.151.
The sign is the same.
The approximate-weight intercept contrast is close: 0.156 does not round to
0.157, and the unrounded gap is 0.00078.
The primary coverage contrast is a match because 0.9449 rounds half up to
0.945, and both sources show weighted multilevel coverage at essentially zero.
The Bernoulli threshold contrast is a mismatch at the thousandth because the
unrounded gap is 0.0027.
The contrast stays negative, the estimator order matches Table C.2, and the
gap is inside two Monte Carlo standard errors on 100 replications.
The random-coefficient slope contrast is a match at three decimals.
Coverage at average cluster size 20 and intraclass correlation .05 is a
mismatch: 0.252 does not round to 0.282.
The raw true-multilevel intercept bias on the primary normal cell is a
mismatch, 0.153 against 0.151, and the contrast with the single-level
estimator is a mismatch under the same 0.001 rule.
The largest non-informative absolute intercept bias is a match at 0.002.
The mean cluster count is a mismatch at the printed cluster and is within one
percent, which is the agreement Table 4.1 supports for a Poisson design.

Small gaps should be read separately from the coverage mismatch.
The raw multilevel bias of 0.153 against 0.151, the one-thousandth contrast
gaps, and the four-cluster difference in sample size are rounding or seed
differences.
They are not replication failures of the headline order.
The Bernoulli gap of 0.003 is a third-decimal mismatch inside the Monte Carlo
band of a 100-replication run, not a sign change.
The weighted single-level random-coefficient intercept, 0.012 against 0.002
with standard error 0.004, is too noisy at 40 replications to call a failure
or a match.
The coverage gap of 0.030, from 0.282 to 0.252, is not a rounding difference.
It is the unresolved cell in the next section.
Both rates are far below the nominal 0.95, so the qualitative report that
these intervals undercover still holds.

\section{What remains unresolved?}

The question is which published results this note does not establish.
The answer is the Bernoulli and random-coefficient cells we did not rerun,
and the coverage gap at average cluster size 20.

Three candidates were open for that coverage gap: a documented finite-sample
multiplier on the sandwich, ordinary Monte Carlo error, and an Mplus standard
error that is not a scalar rescaling of the standard errors we saved.
The saved CR-0 coverage on 984 replications is 0.252, against the printed
0.282 in Table E.2
(\src{simulation/results/coverage_scaling.csv}).
Multiplying the variance, or the standard error, by $G/(G-1)$, by
$N/(N-1)$, or by $G/(G-1)\times(N-1)/(N-3)$ moves the coverage from 0.252 to
0.253.
That is about 0.001, not the gap of 0.030.
The Monte Carlo standard error of our coverage is 0.014 (0.0138 before
rounding to three decimals).
A published rate of 0.282 on 1,000 draws has a binomial standard error of
about 0.014 as well.
The gap is about two of those standard errors, so Monte Carlo error can
account for part of it and cannot show which standard-error definition
produced 0.282.
The largest remaining cause is an Mplus 7 MLR detail that the scalar factors
do not reproduce.
The artifact that would settle it is a Mplus 7 output for the informative
normal cell at average cluster size 20 and intraclass correlation .05, ideally
on the same replicates.
The bias summary for that cell retains 965 replications
(\src{simulation/results/summary.csv}); the scaling file retains 984
replications that have a saved standard error.
Both coverages, 0.248 in the bias file and 0.252 in the scaling file, are far
from 0.282 and far from 0.95.

The other open cells are a different gap.
Bernoulli settings other than intraclass correlation .05 at average cluster
sizes 5 and 20 were not rerun, and the two cells that were rerun used 100
replications rather than 1,000.
Random-coefficient settings other than the informative cell at intraclass
correlation .05 and average cluster size 5 were not rerun, and that cell used
40 replications.
The attempts on those designs are the cells just reported.
They reproduce the headline contrasts and do not substitute for the rest of
the 208 settings.
The remaining cause, ranked ahead of a change to the estimator, is that those
cells were not executed.
The artifact that would settle them is the same R script at the other
intraclass correlations, cluster sizes, and sampling designs, at a replication
count closer to 1,000.
Mplus output on those cells would in addition separate sandwich scaling from
Monte Carlo error.
No Mplus input file was found. The search is recorded in
\src{notes/mplus-search.md}.

\section{Where are the code and the source?}

The question is whether a reader can obtain her programs and our results.
The answer is that our result files are in this repository and her Mplus
inputs were not found.
The dated ledger is \src{notes/availability.md}, searched on 2 November 2026
(UTC).
The Mplus search is \src{notes/mplus-search.md}, with the same date.

The dissertation PDF was recovered from the Internet Archive after the
DigitalCommons download returned an access challenge.
The copy used here is \src{source/koziol2015.pdf}.
The title page matches Natalie A.\ Koziol, May 2015, with adviser James A.\
Bovaird.
Appendix A prints part of the R generator.
The dissertation names Mplus 7 for estimation and R with MASS for generation.
It names no program archive, no OSF project, and no syntax file.
The Figshare supplement of the 2017 article, file
\texttt{hmbr\_a\_1292115\_sm3492.docx} under
\url{https://doi.org/10.6084/m9.figshare.4743673}, repeats simulation tables and
does not contain Mplus or R programs.
Searches of GitHub, the OSF API, Zenodo, the CYFS biography and CV, two later
OSF preregistrations, and the public ORCID record did not produce the
simulation inputs.
The estimation programs were not obtained.
The generator fragment in Appendix A was.

\section{What should happen next?}

The question is which artifact to obtain first, and which check to run on it.
The answer is a Mplus 7 output for the informative normal cell at average
cluster size 20 and intraclass correlation .05.

The check is a comparison of that output's intercept standard errors, and of
the coverage they imply, with our coverage 0.252 and the printed coverage
0.282.
If the output is computed on our replicates, the comparison isolates the
sandwich.
If it is computed on a new Mplus Monte Carlo, the comparison still shows
whether Mplus 7 lands near 0.282 under the options on page 76.
The same file would also show whether \texttt{WTSCALE=CLUSTER} matches the
Method 2 scaling used here.
The next simulation check, which does not require Mplus, is to run the
remaining Bernoulli and random-coefficient cells at the probit
pseudo-likelihood and the cluster-aligned random slope already used for the
cells in this note.

\section{Compliance ledger}

The question is whether the replication requirements that apply to this design
were followed.
The answer is that they were followed, with one departure: the preflight
engine is for structural equation models, and these models are not structural
equation models.

\begin{table}[H]
\centering
\footnotesize
\renewcommand{\arraystretch}{1.3}
\caption{Compliance ledger.
A departure carries its justification in the evidence cell.}
\label{tab:ledger}
\begin{tabular}{@{}p{1.7in}p{0.95in}p{3.3in}@{}}
\toprule
Requirement & Status & Evidence \\
\midrule
Source identified and cited & followed &
Natalie A.\ Koziol (2015), University of Nebraska--Lincoln, chair James A.\ Bovaird.
DigitalCommons URL in the introduction. \\
Verdict before the design & followed &
Abstract and Section 1 open with the replication verdict and the three contrasts. \\
Effect magnitude on page 1 & followed &
Table~\ref{tab:headline}, with tolerance 0.001 and the point-estimate note. \\
Published numbers separated from ours & followed &
Table captions name Tables C.2, C.4, E.2, and 4.1 as the source of published values. \\
Method citations & followed &
Weighting, sandwich variance, pseudo-likelihood, and each package that was discussed are cited in the text. \\
Engine, lavaan / OpenMx & departed &
Departure justification: the preflight expected lavaan and OpenMx.
These models are not structural equation models, so that pair was not used to fit them.
The weighted multilevel estimator was checked by WeMix 4.0.3 against the analytic pseudo-likelihood, and the two agree on the production replicate in \src{simulation/results/wemix_crosscheck_cell1_rep1.csv}. \\
Availability ledger & followed &
\src{notes/availability.md} and \src{notes/mplus-search.md}, dated 2 November 2026. \\
\bottomrule
\end{tabular}
\end{table}

\appendix
\section{Artifacts}
\label{sec:artifacts}

The question is which file holds each result, and what regenerates it.
The answer is Table~\ref{tab:artifacts}.
Links point at the \texttt{master} copy of this repository and are navigation
aids.
The numbers in the note were read from the result files named here.
The invocation, the working directory, and the seeds are recorded in
\src{simulation/results/RUN.txt}.
That file is the regeneration record.
This appendix does not restate it as a second command.

\begin{table}[H]
\centering
\footnotesize
\renewcommand{\arraystretch}{1.25}
\caption{Files that hold the results in this note.
Each link is the file on \texttt{master}.}
\label{tab:artifacts}
\begin{tabular}{@{}>{\raggedright\arraybackslash}p{1.65in}>{\raggedright\arraybackslash}p{2.7in}>{\raggedright\arraybackslash}p{2.15in}@{}}
\toprule
Result & File & Regeneration record \\
\midrule
Normal-grid biases and coverages &
\src{simulation/results/summary.csv} &
\src{simulation/results/RUN.txt}, generator \src{simulation/generate.R} \\
Primary-cell replicates &
\src{simulation/results/cell_01.csv} &
\src{simulation/results/RUN.txt} \\
WeMix cross-check, one replicate &
\src{simulation/results/wemix_crosscheck_cell1_rep1.csv} &
\src{simulation/results/RUN.txt} \\
Bernoulli cell, cluster size 5 &
\src{simulation/results/bern_inf_05_5.csv} &
\src{simulation/results/RUN.txt} \\
Bernoulli cell, cluster size 20 &
\src{simulation/results/bern_inf_05_20.csv} &
\src{simulation/results/RUN.txt} \\
Random-coefficient cell &
\src{simulation/results/rc_inf_05_5.csv} &
\src{simulation/results/RUN.txt} \\
Coverage scaling at cluster size 20 &
\src{simulation/results/coverage_scaling.csv} &
\src{simulation/results/RUN.txt} \\
Dissertation PDF &
\src{source/koziol2015.pdf} &
Obtained from the Internet Archive; see \src{notes/availability.md} \\
Program search &
\src{notes/mplus-search.md} &
Dated ledger, 2 November 2026 \\
Code and supplement ledger &
\src{notes/availability.md} &
Dated ledger, 2 November 2026 \\
Bibliography source &
\src{paper/references.bib} &
Edited with the note \\
\bottomrule
\end{tabular}
\end{table}

\bibliographystyle{plainnat}
\bibliography{references}

\end{document}
